\section{Conclusion}
\label{sec:conclusion}

This research investigated to what extent a supervised object detection model and a multi modal foundation model differ in detecting and identifying predefined inspection relevant categories within housing inspection images. To answer this research question, YOLOv8 and LLaVA 1.6 were evaluated under identical experimental conditions using a harmonized dataset, shared inspection categories, and a common evaluation framework.

The results demonstrate that YOLOv8 consistently outperformed LLaVA 1.6 on the category level classification task. Task specific supervision combined with bounding box annotations enabled more accurate and consistent defect detection than prompt based zero shot reasoning. Although prompt engineering improved the performance of LLaVA, the performance gap with YOLOv8 remained substantial.

Regarding robustness, both approaches were influenced by variations in lighting conditions, viewing angles, image quality, and category characteristics. Nevertheless, YOLOv8 maintained superior performance across the evaluated inspection categories. At the same time, the property level evaluation demonstrated that LLaVA remained capable of correctly identifying overall inspection outcomes after prompt refinement, indicating that foundation models may provide value for higher level inspection reasoning despite weaker category level performance.

This study contributes to the existing literature by providing a controlled comparison between supervised object detection and a multi modal foundation model under identical housing inspection conditions. The results indicate that supervised object detection currently remains the most suitable approach for category level defect detection and localization, whereas vision language models demonstrate greater potential for property level reasoning, inspection reporting, and decision support. Rather than viewing these approaches as competing alternatives, the findings suggest that future housing inspection systems may benefit from combining the complementary strengths of both paradigms.

These conclusions should be interpreted in light of several limitations. The experiments were conducted using one supervised detector, one vision language model, and a dataset consisting of harmonized public datasets combined with images from a single housing corporation. Consequently, the findings are most applicable to similar housing inspection settings and should be validated using additional datasets, organizations, and model architectures.

Future research should investigate hybrid inspection systems that combine accurate object localization with multi modal reasoning. Evaluating newer vision language models, larger housing inspection datasets, and additional techniques inspired by medical imaging, such as improved annotation strategies and methods for detecting subtle visual abnormalities, represents a promising direction for advancing automated housing inspection.